\subsection*{同一次购买可以有多次通信尝试}

把一次 HTTP 尝试编号为 R1，超时后的第二次尝试编号为 R2。两者携带同一 K1。R1 在入口服务留下访问日志，R2 也留下访问日志；即使服务最终只产生一笔购买，运维人员仍需知道两个请求分别到过哪里、等待多久以及返回什么。请求 ID 因而按尝试变化。

分布式追踪用于关联一次调用经过的服务和子调用。一次购买尝试可能经过网关、购买服务和报价服务，各段都有自己的 span。span 是一次追踪中的工作片段，记录一个处理步骤的起止和关联关系；trace 则把这些片段组织成一次可查询的调用经过。重试可以创建新 trace，也可以在已有 trace 下创建新的尝试 span，具体取决于追踪工具。追踪数据可能采样、丢失或按不同期限删除，其 ID 的生成规则也服务于诊断。无论采用哪种组织方式，购买服务都应从明确的业务字段读取 K1。

订单号 O1 标识已形成的订单。它可以在购买事务中才生成，并出现在成功响应里。第一次响应丢失时，客户端可能尚不知道 O1，因此无法要求客户端仅靠 O1 恢复首次购买。幂等键 K1 在首次发送前存在，服务把它关联到 P1；P1 再关联订单 O1 与扣款 D1。返回 O1 后，后续订单查询可以使用 O1，跨请求识别仍沿用 K1。

\begin{figure}[htbp]
\centering
\begin{tikzpicture}[x=1cm,y=1cm]
\node[booknode,text width=3.5cm] (r1) at (2,0) {R1、Trace A\\首次通信与诊断};
\node[booknode,text width=3.5cm] (r2) at (10.4,0) {R2、Trace B\\恢复通信与诊断};
\node[booknode,text width=4cm] (k) at (6.2,-1.6) {K1：购买操作身份\\随两次请求保持稳定};
\node[booknode,text width=4cm] (p) at (6.2,-3.3) {P1：服务保存的购买};
\node[booknode,text width=3.5cm] (o) at (2.4,-4.9) {O1：成交订单};
\node[booknode,text width=3.5cm] (d) at (10,-4.9) {D1：扣款流水};
\draw[bookarrow] (r1)--(k);
\draw[bookarrow] (r2)--(k);
\draw[bookarrow] (k)--(p);
\draw[bookarrow] (p)--(o);
\draw[bookarrow] (p)--(d);
\end{tikzpicture}
\caption{通信 ID 分别保存，稳定身份共同关联到原购买及其订单和扣款。}
\label{fig:ch03-identifiers}
\end{figure}


\begin{table}[htbp]
\centering\small
\begin{tabularx}{\linewidth}{@{}p{3.1cm}p{3.1cm}X@{}}
\toprule
标识 & 常见生成方 & 作用与生命周期\\
\midrule
请求 ID R1 & 客户端或入口服务 & 标识一次 HTTP 尝试，随访问日志保存\\
Trace ID、Span ID & 追踪系统的参与者 & 关联诊断调用，按追踪保留策略保存\\
幂等键 K1 & 本章由客户端生成 & 标识一次购买操作，覆盖约定的恢复期限\\
购买编号 P1 & 购买服务 & 标识业务处理记录，随业务事实保存\\
订单号 O1 & 订单创建方 & 标识成交订单，随订单保存\\
扣款号 D1 & 扣款记录创建方 & 标识一次资金变动，引用 P1 和 O1\\
\bottomrule
\end{tabularx}
\caption{不同标识分别对应通信、诊断与业务记录。}
\end{table}

这些值可以在日志中一起记录，便于沿 R2 找到 K1，再沿 K1 找到 P1 与 O1。诊断关联和业务决策仍各自采用指定字段。若代理为每次重试自动换一个请求 ID，并顺手用它覆盖幂等键，服务将连续收到新操作；若把同一 trace 内所有购买都使用的 trace ID 当作幂等键，又可能错误合并多笔购买。

可以用上一节的划组规则检查一个现有 ID 是否适用。请求 ID 把 R1、R2 分开，追踪 ID 有时又把多项子操作放在一起；它们划分的请求集合都可能与购买授权不一致。一个 ID 是否足够长、是否全局少见，与它是否标识正确对象，是两项独立的问题。只有生成、传播、保存和解释规则都对应到同一次业务操作，服务才能用它恢复原购买。

\subsection*{早期 RPC 怎样表示同一个调用}

第二章介绍的 Birrell 与 Nelson RPC 系统需要辨认结果属于哪个调用，也需要识别通信层重传的调用包。他们把调用身份组织为调用机器、调用进程和调用序号\cite{rpc}。

其中机器与进程共同构成 activity。activity 在这套协议中表示同一机器上的一个调用活动，是接收方保存最近序号的分组单位。协议要求每个 activity 同时至多有一个未完成的远程调用，前一个结果返回后才发起下一个；调用序号沿该 activity 单调增加，即后发起的新调用序号始终大于先前调用序号。接收方保存已看到的序号，判断一个包是新调用还是旧调用的重传。调用编号也随结果返回，使调用者能把晚到结果对应到正确调用。

这个方案里的重传沿用原调用序号。若每次发包都递增序号，接收方得到的就是一串新调用。它还要处理机器重启：重启后的序号若重新落入旧范围，新调用可能被当成重复调用丢掉。方案中因而引入区分调用机器不同运行代次的 conversation identifier，即随机器每次重新启动而换用的会话代次标识，并随调用发送。身份的唯一范围于是包含运行代次，也包含进程和调用序号。

这些前提解释了为什么不能直接给购买客户端只保存一个最近序号。如果一个客户端允许两项购买同时进行，序号 42 的响应可能先于 41 返回。只看到 42 就删除 41 的结果，随后 41 重试时便会失去恢复信息。早期协议每个 activity 至多一个未完成调用的限制，承担了简化状态保存的作用。购买服务允许并发操作时，就需要为仍可能恢复的各项操作分别保留身份与结果。

\begin{figure}[htbp]
\centering
\begin{tikzpicture}[x=1cm,y=1cm]
\node[booknote] at (1.4,.6) {客户端};
\node[booknote] at (11,.6) {服务};
\draw[draw=Ink!30,dashed] (1.4,.2)--(1.4,-3.7);
\draw[draw=Ink!30,dashed] (11,.2)--(11,-3.7);
\draw[bookarrow] (1.4,-.3)--node[above,font=\footnotesize]{发起 41，尚未收到结果}(11,-.7);
\draw[bookarrow] (1.4,-1.2)--node[above,font=\footnotesize]{并发发起 42}(11,-1.6);
\draw[bookarrow] (11,-2)--node[above,font=\footnotesize]{42 的结果先返回}(1.4,-2.4);
\draw[bookalert] (1.4,-2.9)--node[above,font=\footnotesize]{41 仍需恢复，继续携带 41}(11,-3.3);
\node[booknote,text width=12cm] at (6.2,-4.2)
 {已经看到 42，并不能推出 41 已结束；允许并发调用时，保存最近序号不足以覆盖全部未完成操作};
\end{tikzpicture}
\caption{原协议的单个未完成调用前提，不能直接搬到允许并发购买的客户端。}
\label{fig:ch03-sequences}
\end{figure}


这里的关键前提是，一个更大的序号是否足以证明所有更小序号都已结束。如果客户端严格依次发起调用，前一个返回后才发起后一个，接收新调用就能提供一部分先前调用已经结束的证据。允许并发发起后，这项推断失效，41 和 42 可以同时在途。序号仍能区分操作，状态保存方式却要与并发模型一起调整。购买采用独立键时，也需要明确保留哪些尚未确认的结果。

\subsection*{RIFL 与 Raft 把编号接到结果上}

RIFL 将 RPC 身份分成 64 位客户端标识和 64 位客户端序号，同时要求客户端记录该身份，后续重试继续携带它。服务器用这个组合识别新请求、执行中、已完成和已过时的调用；已完成者返回原结果\cite{rifl}。

客户端标识需要系统范围的分配机制，RIFL 采用租约管理。租约是一项有截止时间的资格：客户端定期续期，服务器据其是否仍有效决定能否受理该客户端的旧调用。超期后，服务可按协议回收相关记录，同时拒绝继续沿用旧资格的调用。举一个例子，单独的序号 17 可以在许多客户端上出现，客户端标识使这些 17 分属不同的空间；租约过期后，旧身份也有明确的处理边界。相关协议还处理完成记录的持久性、重试到达的服务器如何找到记录，以及记录的回收。这几项工作分别对应身份以外的故障条件。

Raft 让各副本按共同确定的日志顺序应用命令，领导者负责接收新命令，其他服务器复制日志。这里的提交表示命令已经达到协议要求、可以进入共同应用的序列。Ongaro 与 Ousterhout 的扩展版论文指出，领导者提交命令后、回答客户端前崩溃，客户端向新领导者重试，仍可能把同一命令再次加入日志。其客户端交互方案让客户端为命令分配序号，状态机保存每个客户端已处理的序号及关联响应，发现重试便返回原响应\cite{raft}。

Ongaro 在此基础上还进一步讨论了同一客户端并发调用：保存内容扩展为多组序号与响应，并由客户端报告仍未收到响应的最低序号，帮助服务器回收更早的结果。其会话方案同时涉及注册与过期\cite{ongaro}。本章采用独立购买键，便于逐笔说明输入和账户范围；它与客户端序号方案都要求重试保留身份，如何保存多项未完成操作则取决于各自的并发模型。

\begin{figure}[htbp]
\centering
\begin{tikzpicture}[x=1cm,y=1cm]
\node[booknode,text width=3.4cm] (c) at (1.9,0) {调用方恢复记录\\保留原身份 S};
\node[booknode,text width=3.4cm] (s) at (6.4,0) {恢复请求携带 S\\服务按 S 查找};
\node[booknode,text width=3.4cm] (r) at (10.9,0) {完成记录\\S、业务结果、关联 P1};
\draw[bookarrow] (c)--(s);
\draw[bookarrow] (s)--(r);
\node[booknode,text width=5.5cm] (e) at (10,-2) {与业务变化共同提交\\故障后仍能读取同一记录};
\draw[bookarrow] (r)--(e);
\draw[bookarrow] (e.west)--++(-5,0)-|node[pos=.55,below,font=\footnotesize]{原结果沿身份返回}(c.south);
\node[booknote,text width=12cm] at (6.4,-3.2)
 {身份保存、结果保存和恢复查找共同构成这条路径\\任一段丢失，都需要相应的恢复规则};
\end{tikzpicture}
\caption{稳定身份把重试指向原结果，持久性与查找路径仍须由系统兑现。}
\label{fig:ch03-recovery-chain}
\end{figure}


这些工作把几个经常混在一起的条件分开了。调用方要能恢复原身份；服务处理重试时要能找到同一操作；业务变化成立时要有相应完成记录；该记录在承诺覆盖的故障后仍可恢复。标识相等只能回答该找哪项操作，找到以后返回什么、怎样避免再次生效，还要由保存与执行规则回答。把一个稳定序号放进共识日志，仍须在应用命令时识别它并保存原响应。
\FloatBarrier
